% ID: FIS-OTT-RIF-002
% Titolo: Velocita della luce in un mezzo trasparente
% Disciplina: Fisica
% Area: Ottica
% Argomento: Rifrazione
% Sottoargomento: Legge di Snell
% Classe: Seconda liceo scientifico
% Difficolta: 3
% Tipo: problema_numerico
% Risultato: v circa 2,46 * 10^8 m/s
% Tempo_stimato: 8 min
% Competenze: applicazione della legge di Snell, calcolo dell'indice di rifrazione, modellizzazione
% Prerequisiti: legge di Snell, indice di rifrazione, funzioni goniometriche
% Tag: fisica, ottica, rifrazione, legge_snell, velocita_luce
% Fonte: verifica_fisica_ottica_anonimizzata
% Autore: Riccardo Carli
% Licenza: CC-BY-SA-4.0

\begin{esercizio}
Un fascio di luce si propaga nell'aria e incide su un materiale trasparente.
Gli angoli di incidenza e di rifrazione sono, rispettivamente,
\(\theta_i=63^\circ\) e \(\theta_r=47^\circ\).

Calcola la velocita della luce nel mezzo trasparente.
\end{esercizio}

\begin{soluzione}
Applichiamo la legge di Snell:
\[
n_1\sin\theta_i=n_2\sin\theta_r.
\]
Per l'aria possiamo assumere \(n_1\simeq 1\). Allora:
\[
n_2=\frac{\sin\theta_i}{\sin\theta_r}
=\frac{\sin 63^\circ}{\sin 47^\circ}
\simeq 1{,}22.
\]

La velocita della luce nel mezzo e:
\[
v=\frac{c}{n_2}.
\]

\[
v=\frac{3{,}00\cdot 10^8}{1{,}22}
\simeq 2{,}46\cdot 10^8\,\mathrm{m/s}.
\]
\end{soluzione}
